	摘要是整篇論文的濃縮總結，需精簡扼要，涵蓋論文核心內容，邏輯連貫地總結各章節要點，吸引讀者興趣。建議控制在一頁內，約400–600字，避免使用圖表。
    
\textbf{內容要素：}
\begin{itemize}
    \item 研究動機與目標：闡述研究背景、核心問題及主要目標。
    \item 研究範圍：明確界定研究的課題與範疇。
    \item 研究方法：概述方法、技術或工具，並簡述其合理性。
    \item 主要結果：突出核心發現或數據。
    \item 結果比較與意義：將結果與研究目標及文獻比較，凸顯貢獻。
    \item 未來影響與建議：說明研究對學術或實務的助益，提出未來方向。
\end{itemize}

\textbf{撰寫建議：}
\begin{itemize}
    \item 使用精煉語言，確保非領域專家也能理解。
    \item 完成論文其他部分後再提煉摘要，確保內容全面。
    \item 使用純文字格式（無HTML或特殊格式），便於圖書館建檔與網路檢索。
    \item 附上5–7個關鍵詞，提升論文檢索率。
\end{itemize}

\medskip
\textbf{關鍵字：}論文格式、圖書館建檔、範例、提煉、圖書館建檔
